\subsection{ویژگی‌های استخراج‌شده از هر سیگنال}
\label{app:features}

در این پژوهش، برای هر کانال از سیگنال‌های شتاب‌سنج و ژیروسکوپ، مجموعه‌ای از ویژگی‌های حوزه زمان استخراج شد. اگر سیگنال ورودی به صورت
$x = \{x_1, x_2, \ldots, x_N\}$
در نظر گرفته شود، که در آن $N$ تعداد نمونه‌های سیگنال است، ویژگی‌های استخراج‌شده مطابق جدول
\ref{tab:features}
تعریف شدند. در 
\autoref{app:featuresformula}
، نحوه محاسبه هر یک از این ویژگی‌ها نیز ارائه شده است.

\begin{table}[H]
	\centering
	\caption{ویژگی‌های استخراج‌شده از هر سیگنال}
	\label{tab:features}
	\makebox[\textwidth][c]{%
		\begin{tabular}{p{0.05\textwidth} p{0.30\textwidth} p{0.30\textwidth} p{0.40\textwidth}}
			\toprule
			ردیف & نام فارسی & نام انگلیسی & توضیح\\
			\midrule
			1  & تعداد نمونه‌ها & \lr{Count} & تعداد کل نمونه‌های موجود در سیگنال\\
			2  & میانگین & \lr{Mean} & مقدار متوسط سیگنال\\
			3  & انحراف معیار & \lr{Standard Deviation} & میزان پراکندگی نمونه‌ها نسبت به میانگین\\
			4  & واریانس & \lr{Variance} & شاخص پراکندگی سیگنال\\
			5  & کمینه & \lr{Minimum} & کمترین مقدار ثبت‌شده در سیگنال\\
			6  & بیشینه & \lr{Maximum} & بیشترین مقدار ثبت‌شده در سیگنال\\
			7  & دامنه تغییرات & \lr{Range} & اختلاف بین بیشینه و کمینه سیگنال\\
			8  & میانه & \lr{Median} & مقدار مرکزی سیگنال پس از مرتب‌سازی نمونه‌ها\\
			9  & چارک اول & \lr{First Quartile} & صدک ۲۵ام سیگنال\\
			10 & چارک سوم & \lr{Third Quartile} & صدک ۷۵ام سیگنال\\
			11 & دامنه بین‌چارکی & \lr{Interquartile Range} & اختلاف بین چارک سوم و چارک اول\\
			12 & میانگین قدر مطلق انحرافات & \lr{Mean Absolute Deviation} & میانگین فاصله قدر مطلق نمونه‌ها از میانگین\\
			13 & میانگین قدر مطلق & \lr{Mean Absolute Value} & میانگین اندازه مطلق مقادیر سیگنال\\
			14 & ریشه میانگین مربعات & \lr{Root Mean Square} & شاخصی از شدت کلی سیگنال\\
			15 & انرژی سیگنال & \lr{Energy} & مجموع توان نمونه‌های سیگنال\\
			16 & ضریب تغییرات & \lr{Coefficient of Variation} & نسبت انحراف معیار به میانگین\\
			17 & چولگی & \lr{Skewness} & شاخص عدم تقارن توزیع مقادیر سیگنال\\
			18 & کشیدگی & \lr{Kurtosis} & شاخص میزان تیزی یا پهنی توزیع سیگنال\\
			19 & آنتروپی & \lr{Entropy} & شاخصی از میزان بی‌نظمی یا پراکندگی توزیع سیگنال\\
			20 & تعداد عبور از صفر & \lr{Zero Crossings} & تعداد دفعات تغییر علامت سیگنال\\
			21 & میانگین تفاضل نمونه‌های متوالی & \lr{Mean of First Differences} & میانگین تغییرات بین نمونه‌های متوالی\\
			22 & انحراف معیار تفاضل نمونه‌های متوالی & \lr{Standard Deviation of First Differences} & میزان پراکندگی تغییرات متوالی سیگنال\\
			23 & خودهمبستگی با تأخیر یک & \lr{Lag-1 Autocorrelation} & میزان وابستگی هر نمونه به نمونه بعدی\\
			\bottomrule
		\end{tabular}
	}
\end{table}

\newpage
\subsubsection{فرمول محاسبه ویژگی‌های استخراج‌شده}
\label{app:featuresformula}

در روابط زیر، سیگنال ورودی به صورت $x = \{x_1, x_2, \ldots, x_N\}$ در نظر گرفته شده است. همچنین $\mu$ میانگین سیگنال، $\sigma$ انحراف معیار، $Q_1$ چارک اول، $Q_3$ چارک سوم و $\Delta x_i = x_{i+1}-x_i$ تفاضل نمونه‌های متوالی را نشان می‌دهد.
\begin{table}[H]
	\centering
	\caption{ویژگی‌های استخراج‌شده از هر سیگنال و نحوه محاسبه آن‌ها}
	\label{tab:featuresformula}
	\makebox[\textwidth][c]{%
	\begin{tabular}{p{0.30\textwidth} p{0.50\textwidth}}
	
		\toprule
		\textbf{ویژگی} & \textbf{رابطه}\\
		\midrule
		
		تعداد نمونه‌ها
		&
		$N$
		\\
		
		میانگین
		&
		$\mu = \frac{1}{N}\sum_{i=1}^{N} x_i$
		\\
		
		انحراف معیار
		&
		$\sigma = \sqrt{\frac{1}{N-1}\sum_{i=1}^{N}(x_i-\mu)^2}$
		\\
		
		واریانس
		&
		$s^2 = \frac{1}{N-1}\sum_{i=1}^{N}(x_i-\mu)^2$
		\\
		
		کمینه
		&
		$\min(x) = \min\{x_1,\ldots,x_N\}$
		\\
		
		بیشینه
		&
		$\max(x) = \max\{x_1,\ldots,x_N\}$
		\\
		
		دامنه تغییرات
		&
		$Range = \max(x)-\min(x)$
		\\
		
		میانه
		&
		$\tilde{x} = median(x)$
		\\
		
		چارک اول
		&
		$Q_1 = P_{25}(x)$
		\\
		
		چارک سوم
		&
		$Q_3 = P_{75}(x)$
		\\
		
		دامنه بین چارکی
		&
		$IQR = Q_3 - Q_1$
		\\
		
		میانگین قدر مطلق انحرافات
		&
		$MAD = \frac{1}{N}\sum_{i=1}^{N}|x_i-\mu|$
		\\
		
		میانگین قدر مطلق
		&
		$MeanAbs = \frac{1}{N}\sum_{i=1}^{N}|x_i|$
		\\
		
		ریشه میانگین مربعات
		&
		$RMS = \sqrt{\frac{1}{N}\sum_{i=1}^{N}x_i^2}$
		\\
		
		انرژی سیگنال
		&
		$Energy = \sum_{i=1}^{N}x_i^2$
		\\
		
		ضریب تغییرات
		&
		$CV = \frac{\sigma}{|\mu|}$
		\\
		
		چولگی
		&
		$Skewness = \frac{1}{N}\sum_{i=1}^{N}\left(\frac{x_i-\mu}{\sigma}\right)^3$
		\\
		
		کشیدگی
		&
		$Kurtosis = \frac{1}{N}\sum_{i=1}^{N}\left(\frac{x_i-\mu}{\sigma}\right)^4$
		\\
		
		آنتروپی
		&
		$Entropy = -\sum_{j=1}^{B}p_j \log(p_j)$
		\\
		
		تعداد عبور از صفر
		&
		$ZC = \sum_{i=1}^{N-1} I(x_i x_{i+1}<0)$
		\\
		
		میانگین تغییرات متوالی
		&
		$DiffMean = \frac{1}{N-1}\sum_{i=1}^{N-1}(x_{i+1}-x_i)$
		\\
		
		انحراف معیار تغییرات متوالی
		&
		$DiffStd = std(\Delta x)$
		\\
		
		خودهمبستگی با تأخیر یک
		&
		$Autocorr_1 = 
		\frac{\sum_{i=1}^{N-1}(x_i-\mu)(x_{i+1}-\mu)}
		{\sum_{i=1}^{N}(x_i-\mu)^2}$
		\\
		
		\bottomrule
	\end{tabular}
}
\end{table}